% ================= 6. TESTING AND VALIDATION PLAN =================
\section{Testing and Validation Plan}
\label{sec:testing}

This plan covers system-level validation. Parameter-identification tests that
feed the model directly (swing test, current-step identification) are described
with the model they validate in Section~\ref{sec:paramid}.

The plan is written before the design point is closed. Its methods, metrics and
orderings are therefore fixed here, while the numeric thresholds that depend on
the cube mass, edge length, brake torque or wheel geometry are carried as
\textcolor{red}{TBD} and populated by the closure procedure of
Section~\ref{sec:massbudget}. This asymmetry is deliberate: a validation plan
whose thresholds are pending is a plan awaiting data, whereas one whose methods
are pending is not yet a plan. Where a threshold is already fixed elsewhere in
this document --- a gate criterion of Table~\ref{tab:gates}, or a device limit
from Appendix~\ref{app:bom} --- it is cited here rather than restated as a new
requirement.

\subsection{Validation Strategy and Stage Gates}
\label{sec:test_strategy}

Validation proceeds in ten stages, ordered so that each one removes the
uncertainty that would otherwise invalidate the next. Three orderings in
particular are not matters of convenience. Feasibility is verified before
fabrication (V1 before V4), because a wheel that cannot meet its inertia target
is cheaper to discover in a spreadsheet than in filament. The control law is
exercised in simulation before it is exercised on hardware (V3 before V7),
because a divergent gain set is diagnosed safely in simulation and destructively
on a rig. Plant identification precedes control validation (V6 before V7),
because gains derived from assumed parameters are not the gains that will be
flown.

The governing rule is that no stage is entered until the exit criterion of the
preceding stage is met. Where a stage fails, the response is to return to the
stage that produced the offending assumption rather than to proceed with a
known-invalid input; Section~\ref{sec:schedule} records the contingency for the
schedule consequences of doing so.

\begin{table}[h]
\centering
\caption[Validation stages]{Validation stages. Exit gates refer to Table~\ref{tab:gates}; no gate
or date is introduced here that is not already defined there.}
\label{tab:validation_stages}
\small
\setlength{\tabcolsep}{4pt}
\begin{tabular}{c L{3.5cm} L{4.3cm} L{3.4cm} c}
\toprule
\# & Stage & Entry condition & Pass criterion & Exit gate \\
\midrule
V1  & Pre-fabrication feasibility (Sec.~\ref{sec:test_feasibility}) & Candidate design point from Sec.~\ref{sec:massbudget} & Both jump-up routes agree; margin held & G2 \\
V2  & CAD and dimensional verification (Sec.~\ref{sec:test_cad}) & V1 passed; model complete & No interference; mass roll-up within budget & G2 \\
V3  & Simulation-first control (Sec.~\ref{sec:test_sim}) & 2D model of Sec.~\ref{sec:control2d} & Stable under all test cases and perturbation & G2 \\
V4  & Fabrication and mechanical acceptance (Sec.~\ref{sec:test_cad}) & V1, V2 passed & Parts within dimensional tolerance; wheel balanced & pre-G3 \\
V5  & Electrical bring-up (Sec.~\ref{sec:test_electrical}) & V4 parts available & All criteria of Table~\ref{tab:electrical_criteria} & G3, G5 \\
V6  & 1-DoF plant identification (Sec.~\ref{sec:test_paramid}) & V5 single axis live & Four parameters measured to tolerance & G3 \\
V7  & 1-DoF control validation (Sec.~\ref{sec:test_1dof}) & V3, V6 passed & Rig balances $\geq$60 s unassisted & G3 \\
V8  & Edge balance (Sec.~\ref{sec:test_edge}) & V7 passed; cube integrated & $\geq$60 s with disturbance recovery & \textbf{G6} \\
V9  & Corner balance and slew (Sec.~\ref{sec:test_corner}) & V8 passed & Point-contact regulation; slew tracked & G7 \\
V10 & Jump-up (Sec.~\ref{sec:test_jumpup}) & V8 passed (hard kill) & Face-to-edge transition captured & G8 \\
\bottomrule
\end{tabular}
\end{table}

Gate G6 is a hard kill for the jump-up objective: if edge balance is not
demonstrated by 13~August, V10 is abandoned rather than attempted, since a cube
that cannot hold an edge cannot re-stabilise after arriving on one. This is a
consequence of the objective ordering of Section~\ref{sec:objectives}, in which
jump-up is a capability built on balance rather than an alternative to it.

\subsection{Pre-Fabrication Feasibility Verification}
\label{sec:test_feasibility}

The purpose of this stage is to establish, before any part is committed to
fabrication, that the wheel about to be built can deliver the jump-up manoeuvre.
This is a stricter question than whether the wheel meets its inertia target,
because the target is itself the output of a model whose assumptions are not
exact.

\subsubsection{Two Independent Routes, Compared}

Section~\ref{sec:jumpup} contains two derivations of the jump-up requirement that
proceed from different assumptions, and this stage evaluates both rather than
selecting one. The collision relation of Section~\ref{sec:jumpup} treats the
wheel--body interaction as a perfectly inelastic impact and returns the required
wheel speed $\omega_w$ directly from the body inertia and mass distribution. The
torque-limited momentum budget of Equations~\eqref{eq:hw_ideal}
and~\eqref{eq:Iw_target} instead treats the achievable wheel speed as fixed and
returns the required inertia $I_{w,\text{target}}$, while additionally accounting
for a finite brake torque that the collision relation idealises away.

Agreement between the two is the evidence that the sizing is sound; disagreement
is a finding in its own right and is reported rather than reconciled by choosing
the more favourable result. The comparison is recorded in
Table~\ref{tab:jumpup_gate}.

\begin{table}[h]
\centering
\caption[Pre-fabrication feasibility gate]{Pre-fabrication feasibility gate. Pass conditions are stated as
relations because the quantities they relate are outputs of the closure
procedure of Section~\ref{sec:massbudget}, not inputs fixed here.}
\label{tab:jumpup_gate}
\small
\setlength{\tabcolsep}{4pt}
\begin{tabular}{L{4.2cm} L{3.5cm} c L{4.0cm}}
\toprule
Quantity & Source & Value & Pass condition \\
\midrule
Required wheel speed $\omega_w$ & Collision relation, Sec.~\ref{sec:jumpup} & \textcolor{red}{TBD} & $\leq$ achievable speed with margin \\
Achievable wheel speed $\omega_{\max}$ & Drive analysis, Sec.~\ref{sec:jumpup_target} & \textcolor{red}{TBD} & Sustained in service, not no-load \\
Speed margin $\omega_w/\omega_{\max}$ & Derived & \textcolor{red}{TBD} & $<1$ by a stated margin \\
Target inertia $I_{w,\text{target}}$ & Eq.~\eqref{eq:Iw_target} & \textcolor{red}{TBD} & --- \\
Achieved inertia $I_w$ & Wheel design, Sec.~\ref{sec:rw_sizing} & \textcolor{red}{TBD} & $\geq I_{w,\text{target}}$ \\
Route disagreement & Both of the above & \textcolor{red}{TBD} & Reported; explained if large \\
Ballast trim authority & Sec.~\ref{sec:rw_sizing} & \textcolor{red}{TBD} & Non-zero in both directions \\
Nut retention margin & Sec.~\ref{sec:rw_sizing} & \textcolor{red}{TBD} & Retention load carried by a stated feature \\
\bottomrule
\end{tabular}
\end{table}

\subsubsection{Brake-Torque Sensitivity}

The brake-amplification factor $\beta = \tau_b/(\tau_b - \tau_g)$ of
Equation~\eqref{eq:Iw_target} diverges as $\tau_b \to \tau_g$, so the inertia
target is not merely uncertain in $\tau_b$ but unboundedly sensitive to it near
the floor. A single assumed value of $\tau_b$ would therefore carry an unstated
and potentially unbounded error into every downstream number.

This stage consequently requires $\tau_b$ to be reported as a sensitivity sweep
rather than as a point value: $I_{w,\text{target}}$ is tabulated across a range
of candidate brake torques, and the achieved wheel inertia is checked against
each. The output of interest is the lowest brake torque at which the design still
closes, which converts an unmeasured parameter into a stated requirement on the
brake. That requirement is discharged when the brake torque is measured on
hardware in Section~\ref{sec:test_electrical}; until then the sweep, not an
assumption, is what the design rests on.

\subsubsection{Trim Authority}

Because $h_w \propto M L^{3/2}$, a mass overrun raises the inertia requirement
rather than leaving it fixed, and the wheel must therefore retain the ability to
be corrected in both directions after the cube is weighed. The ballast scheme of
Section~\ref{sec:rw_sizing} provides this as a discrete trim, and this stage
records how much of it remains unspent at the selected operating point: a design
sitting only marginally above its target has no authority to absorb a mass
overrun and is recorded as failing this gate even when the nominal inertia check
passes. The trim is reported as the number of ballast elements that may be added
or removed before the target is missed.

\subsection{CAD and Dimensional Verification}
\label{sec:test_cad}

The checks below split into those that run against the model, before anything is
fabricated, and those that run against the physical part afterwards. The driving
dimensions themselves are outputs of the packaging closure of
Section~\ref{sec:massbudget} and are not verified here; what is verified is that
the model realises them consistently and that the fabricated part matches the
model.

\paragraph{Model checks (pre-fabrication).} Interference and clearance
verification on the wheel-to-subframe and inter-module clearances $c_w$ and $c_o$
of Table~\ref{tab:cad_params}, through a full rotation of each wheel and across
all three axes simultaneously; orthogonality of the three motor axes; position of
the centre of mass relative to the geometric centre, since the modelling of
Section~\ref{sec:modeling} assumes a known offset; and the mass roll-up against
Table~\ref{tab:massbudget}, which is the input to the feasibility gate of
Section~\ref{sec:test_feasibility} and must be re-run whenever it changes.

\paragraph{Part checks (post-fabrication).} Dimensional acceptance of the ballast
pockets and their retention floor, since both are structural features rather than
cosmetic ones and a print that is dimensionally out of tolerance at the pocket is
a wheel that cannot be trusted at speed; measured wheel mass against the CAD
prediction; and static balance of the assembled wheel, which is the first
opportunity to detect a ballast placement error before it appears as vibration at
speed.

The modelled-versus-measured comparison is recorded in
Table~\ref{tab:mass_verification} of Appendix~\ref{app:cad_mass} and is not
duplicated here. A discrepancy found at that step is carried back into
$I_{w,\text{target}}$ and absorbed by ballast trim where the authority allows,
per Section~\ref{sec:test_feasibility}; where it does not, the finding returns to
the closure procedure rather than being accepted.

\subsection{Simulation-First Validation}
\label{sec:test_sim}

The control law is validated in simulation on the 2D model of
Section~\ref{sec:control2d} before any closed-loop hardware run. The plant is
constructed from the parameters of Table~\ref{tab:params}, which are estimates
until Section~\ref{sec:test_paramid} measures them.

\paragraph{Test cases.} Three cases are run against the linear-quadratic
regulator baseline: recovery from an initial tilt offset, which establishes the
region of attraction; rejection of an impulse disturbance applied at the body,
which is the simulated counterpart of the disturbance-recovery criterion at G6;
and behaviour under wheel-speed saturation, which exercises the momentum-dumping
path of Section~\ref{sec:control_sw} and is the case most likely to be missed by
a purely linear analysis, since saturation is precisely where the linear model
ceases to describe the plant.

\paragraph{Robustness.} Each case is re-run with the plant parameters perturbed
about their nominal values, by a magnitude \textcolor{red}{TBD} set once the
identification tolerances of Section~\ref{sec:test_paramid} are known. The
purpose is to establish that the gains are not knife-edge: a gain set that
stabilises only the nominal plant is not a result, because the nominal plant is
not the one that will be built.

\paragraph{Direction of validation.} Simulation predicts and hardware corrects.
The measured parameters of Section~\ref{sec:test_paramid} are substituted into
the plant and the gains re-derived before hardware balance is attempted; the
simulation is validated against the measured plant, never the reverse. A
discrepancy between simulated and measured closed-loop behaviour is treated as a
finding about the model, and Section~\ref{sec:test_perf} defines the metrics in
which that comparison is expressed.

\subsection{Electrical Bring-up and Isolation}
\label{sec:test_electrical}

Bring-up is staged so that each power step is taken only once the preceding one
has been shown safe: a current-limited bench supply before the battery, wheels
removed before wheels fitted, and a single axis before three. The rationale is
that each step removes one class of fault while the energy available to it is
still bounded --- the first battery connection is not an occasion to discover a
wiring error. The detailed bring-up and isolation checklists are given in
Appendix~\ref{app:electrical} and are not restated here; what follows is the set
of numeric criteria that close each step.

\begin{table}[h]
\centering
\caption[Electrical acceptance criteria]{Electrical acceptance criteria. Task references are to the electrical
workstream schedule, which is the source of each figure.}
\label{tab:electrical_criteria}
\small
\setlength{\tabcolsep}{4pt}
\begin{tabular}{L{3.9cm} L{7.4cm} c}
\toprule
Item & Criterion & Task \\
\midrule
Motor incoming inspection & Phase-to-phase resistance within $\pm$20\% of the \qty{194}{\milli\ohm} datasheet figure on all three pairs; driver enumerates at a unique CAN ID & EL-13 \\
Encoder air gap & $\leq$\qty{0.5}{\milli\meter}, recorded as a numeric shim value, chip centred on the magnet; monotonic absolute angle over a full revolution with zero error flags & EL-14 \\
Logic rail & \qty{5.00}{\volt} $\pm$\qty{0.05}{\volt} at no load; $\geq$\qty{4.75}{\volt} at \qty{1.5}{\ampere}; ripple and buck case temperature recorded & EL-15 \\
CAN bus, rig & Unpowered CANH--CANL \qty{60}{\ohm} $\pm$\qty{3}{\ohm}; all nodes enumerate at \qty{5}{\mega\bit\per\second} & EL-16 \\
CAN bus, integrated & \qty{60}{\ohm} $\pm$\qty{3}{\ohm} with termination at the two physical ends only; zero error frames over a 30-minute soak with all three drivers commanded & EL-24 \\
Bus transient & Bus holds $\geq$\qty{20}{\volt} at the three-axis current transient, measured at the board rather than the battery terminals; no fuse opening, no driver undervoltage fault & EL-25 \\
Telemetry rail disturbance & Logic rail sag $<$\qty{100}{\milli\volt} during a \qty{300}{\milli\ampere} transmit burst; packet loss stated as a percentage & EL-22 \\
Brake servo & Full commanded range on \qty{50}{\hertz} PWM; logic rail disturbance $<$\qty{100}{\milli\volt} during a servo step; flyback path verified & EL-31 \\
Integrated power-on & All three axes addressable on battery power, all encoders within the gap spec after assembly, zero CAN error frames during a three-axis open-loop spin & EL-27 \\
\bottomrule
\end{tabular}
\end{table}

The brake torque $\tau_b$ assumed in the sensitivity sweep of
Section~\ref{sec:test_feasibility} is measured during this stage. Because the
sweep establishes the lowest brake torque at which the design closes, the
measurement is a pass/fail against that figure rather than an open-ended
characterisation.

\subsection{Sensor Testing}
\label{sec:test_sensors}

Section~\ref{sec:background} identifies sensing, rather than the control law, as
the effect that sets the achievable balancing limit in practical replications of
this system. This subsection is the evidence for that claim, and its results
bound what the control performance of Section~\ref{sec:test_perf} can be expected
to achieve.

\paragraph{Inertial measurement.} Static Allan deviation of the gyroscope over a
logging run long enough to resolve the bias-instability floor, which gives the
drift rate the attitude estimate must be corrected against; accelerometer noise
floor in the same conditions; and the resulting bound on the tilt estimate as a
function of the complementary-filter cross-over frequency.

\paragraph{Encoder.} Air-gap verification to the recorded assembly spec, repeated
after mechanical bolt-up, since the gap can shift when the structure is
assembled --- this is why the re-check appears as a separate step in
Table~\ref{tab:assembly_seq} rather than being folded into the initial fitting.
Linearity is checked over a full revolution against a known reference, and error
flags are monitored over a sustained spin rather than sampled once.

\paragraph{Fusion.} The complementary filter of Section~\ref{sec:control2d} is
validated against a known static tilt and a known dynamic profile, and the
cross-over frequency is selected from the measured noise and drift figures above
rather than assumed. The estimator's behaviour under the accelerometer trust
gating of Appendix~\ref{app:control_algorithms} is exercised deliberately, by applying a
body acceleration large enough to trip the gate, since an estimator whose gating
is never tested is one whose gating is not known to work.

\subsection{1-DoF Plant Identification}
\label{sec:test_paramid}

Four parameters are measured on the single-axis rig and returned to the control
design of Section~\ref{sec:test_sim}. Each is measured by a method chosen so that
the quantity of interest dominates the measurement rather than appearing as a
small correction to something else.

\begin{table}[h]
\centering
\caption[Plant parameters measured on the 1-DoF rig]{Plant parameters measured on the 1-DoF rig. Expected values follow from
the design point and are populated when it closes.}
\label{tab:paramid_methods}
\small
\setlength{\tabcolsep}{4pt}
\begin{tabular}{L{2.9cm} L{6.1cm} L{2.6cm} c}
\toprule
Parameter & Method & Instrument & Value \\
\midrule
Torque constant $K_t$ & Current step against a locked rotor; torque from measured reaction & Driver current telemetry & \textcolor{red}{TBD} \\
Body inertia $J$ & Free-swing test about the pivot; inertia from the measured oscillation period & Encoder or IMU log & \textcolor{red}{TBD} \\
Friction & Wheel coast-down from speed; viscous and Coulomb terms separated by fit & Encoder log & \textcolor{red}{TBD} \\
Loop latency & Command-to-response timestamp through the full sense--compute--actuate path & Loopback on the flight computer & \textcolor{red}{TBD} \\
\bottomrule
\end{tabular}
\end{table}

Tolerances are stated with the measured values, since a parameter without a
tolerance cannot support the robustness sweep of Section~\ref{sec:test_sim} that
consumes it. Loop latency is measured through the complete path rather than
estimated from the \qty{1}{\kilo\hertz} loop rate, because the scheduling and
bus-transport contributions are the ones a rate figure omits.

\subsection{1-DoF Control Validation}
\label{sec:test_1dof}

The single-axis rig is the protected experimental item of
Section~\ref{sec:schedule}, and its balance result is the evidence carried into
the document. Gate G3 requires the rig to balance for at least 60~s unassisted.

Beyond the duration criterion, three quantities are recorded because they
determine whether the result transfers to the integrated cube rather than merely
holding on the rig. Recovery from a defined impulse establishes that the balance
is a stabilised equilibrium and not a fortunate initial condition. Steady-state
wheel-speed drift measures momentum accumulation: a controller that balances
while monotonically spinning up is one that will saturate, and the drift rate
predicts when. Control effort is compared against the continuous and peak current
limits of the driver given in Appendix~\ref{app:bom}, since a balance sustained
only by exceeding the continuous rating is not a balance that survives to the
demonstration.

\subsection{Edge, Corner and Jump-Up Validation}
\label{sec:test_milestones}

The three hardware milestones follow the objective ordering of
Section~\ref{sec:objectives}, each entered only on the passing of the previous.

\subsubsection{Edge Balance}
\label{sec:test_edge}

Regulation about a line-contact equilibrium, and the hard-kill gate for the
jump-up objective. G6 requires balance for at least 60~s with disturbance
recovery. Recorded: sustained duration; the maximum initial tilt from which the
controller recovers, which is the experimental region of attraction and is
compared against the simulated figure from Section~\ref{sec:test_sim}; the
disturbance-rejection bandwidth; and wheel-speed drift as in
Section~\ref{sec:test_1dof}. Numeric targets beyond the gate criterion are
\textcolor{red}{TBD}.

\subsubsection{Corner Balance and Commanded Slew}
\label{sec:test_corner}

Regulation about a point-contact equilibrium on all three axes simultaneously,
which is the case the wheel was sized against and therefore the first test in
which the sizing assumption of Section~\ref{sec:test_feasibility} is exercised
directly. Recorded: sustained duration, per-axis control effort, and cross-axis
coupling, the last being absent from the 2D model and observable only here.

\label{sec:test_slew}
Commanded rotation between equilibria is verified in the same stage: the cube is
commanded through a slew and the tracking error against the reference recorded,
together with the settling behaviour on arrival. Targets \textcolor{red}{TBD}.

\subsubsection{Jump-Up}
\label{sec:test_jumpup}

The face-to-edge transition, captured on video at a frame rate sufficient to
resolve the manoeuvre, per G8.

This test closes the loop opened in Section~\ref{sec:test_feasibility}. The wheel
speed at brake release is measured and compared against the value predicted
before fabrication, and the two routes of Section~\ref{sec:test_feasibility} are
re-evaluated against the measured brake torque. The discrepancy is the
quantitative measure of the inelastic-collision idealisation, and its measurement
is the experimental refinement of $\omega_w$ that Section~\ref{sec:jumpup}
anticipates. Recorded: wheel speed at release, achieved versus predicted; whether
the edge is reached and held; and, on a successful transition, the time to
re-stabilise, which is the quantity that determines whether chained jump-up ---
the walking objective --- is reachable.

\subsection{Control Performance Evaluation}
\label{sec:test_perf}

The metrics used above are defined once here so that simulated and measured
results are expressed in the same terms and remain comparable across stages.

\begin{table}[h]
\centering
\caption[Control performance metrics]{Control performance metrics. Definitions and methods are fixed;
targets follow from the closed design point.}
\label{tab:perf_metrics}
\small
\setlength{\tabcolsep}{4pt}
\begin{tabular}{L{2.8cm} L{4.7cm} L{4.1cm} c}
\toprule
Metric & Definition & Measurement method & Target \\
\midrule
Recovery angle & Largest initial tilt from which the equilibrium is regained & Ramped initial-condition release & \textcolor{red}{TBD} \\
Settling time & Time to enter and remain within a stated band about the equilibrium & Step or impulse response log & \textcolor{red}{TBD} \\
Disturbance-rejection bandwidth & Frequency to which body disturbances are attenuated & Swept or impulsive disturbance & \textcolor{red}{TBD} \\
Slew tracking error & Peak and RMS deviation from the commanded trajectory & Reference versus estimated attitude & \textcolor{red}{TBD} \\
Control effort & RMS and peak wheel current over the run & Driver current telemetry & Within Appendix~\ref{app:bom} limits \\
Momentum drift & Rate of change of mean wheel speed while regulating & Encoder log over a sustained run & \textcolor{red}{TBD} \\
\bottomrule
\end{tabular}
\end{table}

\subsection{Milestone Acceptance Criteria}
\label{sec:test_acceptance}

Table~\ref{tab:acceptance} states the condition under which each objective of
Table~\ref{tab:traceability} is accepted as demonstrated, and the subsection that
produces the evidence. It is the closing element of the traceability chain that
runs from the objectives of Section~\ref{sec:objectives}, through the required
capabilities of Section~\ref{sec:requirements}, to the test that confirms them.

\begin{table}[h]
\centering
\caption{Objective acceptance criteria and the evidence that supports each.}
\label{tab:acceptance}
\small
\setlength{\tabcolsep}{4pt}
\begin{tabular}{L{3.0cm} L{6.4cm} c c}
\toprule
Objective & Acceptance condition & Evidence & Gate \\
\midrule
Edge balance & Balanced $\geq$60 s unassisted with recovery from a defined disturbance & Sec.~\ref{sec:test_edge} & G6 \\
Corner balance & Point-contact regulation sustained on all three axes & Sec.~\ref{sec:test_corner} & G7 \\
Controlled rotation & Commanded slew completed within the stated tracking error & Sec.~\ref{sec:test_slew} & G7 \\
Jump-up & Face-to-edge transition achieved and captured on video & Sec.~\ref{sec:test_jumpup} & G8 \\
Walking (stretch) & Chained jump-up with re-stabilisation between transitions & Sec.~\ref{sec:test_jumpup} & --- \\
\bottomrule
\end{tabular}
\end{table}
